\chapter{Der Node-Workflow}\label{kap:workflow}\index{Node-Workflow}

Viele optische Aufgaben lassen sich nicht mit \emph{einem} Rechenverfahren lösen. Ein
Auflichtmikroskop etwa besteht aus Objektiven und Strahlteilern, die man mit Strahlen rechnet,
einer Probe, an der das Licht gebeugt wird (ein Gitter, ein Schichtstapel, ein Partikel), und
einer Kamera nahe einem Fokus, wo Strahlen versagen und Wellenoptik nötig ist. Der
\emph{Node-Workflow} setzt eine solche Rechnung aus \emph{Blöcken} zusammen: Jeder Block ist ein
optisches Element oder eine Auswertung, Verbindungen zeigen, wohin das Licht läuft, und jeder
Block wird von dem Rechenkern gerechnet, der ihn am besten kann.

Dieses Kapitel erklärt die Idee, den Editor im Fenster \gui{Node workflow}, die Äste, Gruppen und
Varianten, das Prüfen und Rechnen mit dem Plugin \plugin{Workflow}, die Anzeigefenster und das
Speichern der Ergebnisse. Alle Blocktypen mit ihren Anschlüssen und Parametern stehen in
Anhang~\ref{kap:workflow-bloecke}.

% ---------------------------------------------------------------------------
\section{Die Idee}\label{sec:workflow-idee}

\subsection{Blöcke, Ports und Verbindungen}

Ein Workflow besteht aus:

\begin{description}
  \item[Blöcken] (engl.\ \emph{blocks} oder \emph{nodes}) – Lichtquellen, optische Elemente
        (Linsengruppen, Spiegel, Strahlteiler, Gitter, Schichtstapel, ideale Linsen),
        Detektoren bzw.\ Kameras, Umwandlungen, Meta-Blöcke (Parameterstudien, Optimierung)
        und Anzeigeblöcke. Jeder Block hat einen \emph{Typ} (z.\,B.\ \block{rtt.section}),
        eine Kennung (z.\,B.\ \code{objective}), einen Namen und Parameter.
  \item[Ports] – die Stellen, an denen Licht in einen Block hinein- oder aus ihm herausgeht,
        z.\,B.\ \code{front} und \code{back} einer Linse. Ein Port ist eine Referenzfläche des
        Blocks und trägt Licht in \emph{beiden} Richtungen.
  \item[Verbindungen] (engl.\ \emph{connections}) – sie verbinden zwei Ports zweier Blöcke.
        Eine Verbindung hat keine Richtung: Das Licht kann sie hin und zurück durchlaufen
        (etwa zur Probe hin und von ihr zurück).
  \item[Daten-Anschlüsse] (engl.\ \emph{sockets}) und \emph{Links} – gerichtete Verbindungen
        für Daten, die kein Licht sind: das Bild einer Kamera zu einer Anzeige, der Antrieb
        (\emph{drive}) eines Meta-Blocks zu einem Spiegel, die Ergebnisse eines Sweeps zu einem
        Diagramm.
  \item[Ästen] (engl.\ \emph{branches}) – benannte Lichtwege von einer Quelle über eine Folge
        von Blöcken zu einem Detektor. \emph{Genau die Äste werden gerechnet}; eine Verbindung,
        die kein Ast benutzt, trägt kein Licht (Abschnitt~\ref{sec:workflow-aeste}).
\end{description}

\begin{hinweis}
Ältere Workflows bestanden nur aus gerichteten Anschlüssen (Eingänge links, Ausgänge rechts)
ohne Ports und Äste. Sie lassen sich weiter öffnen, prüfen und rechnen. \gui{Workflow from
system} des Raytracers verwandelt einen solchen Workflow in die neue Form mit Ports und Ästen und
übernimmt dabei Ihre Änderungen (Positionen, Parameter, ersetzte und hinzugefügte Blöcke,
Verbindungen, Transporte, eigene Äste). Dieses Kapitel beschreibt die neue Form; wo sich die alte
unterscheidet, ist es vermerkt.
\end{hinweis}

\subsection{Darstellungen: Strahlen, Felder und Spektren}\label{sec:workflow-darstellungen}\index{Darstellung}

Licht kann auf einer Verbindung in verschiedenen \emph{Darstellungen} (engl.\
\emph{representations}) transportiert werden. Jede hat einen Kennbuchstaben:

\begin{referenz}{@{}>{\bfseries}p{0.06\linewidth}>{\raggedright\arraybackslash}p{0.27\linewidth}X@{}}
\toprule
\textbf{} & \textbf{Darstellung} & \textbf{Inhalt} \\
\midrule
\endhead
R & Strahlen & Strahlbündel mit Ort, Richtung und optischer Weglänge (damit auch Phase);
  die schnellste Darstellung. Linsen, Spiegel und Strahlteiler eines optischen Systems
  (\block{rtt.*}) rechnen immer mit Strahlen. \\
H & hybrides Feld & Feld ohne Gitter (Amplitude und Phase mit ihren Ableitungen), gerechnet
  mit \emph{Fieldtracing}; für Fokusse, Beugung und Interferenz. \\
S & Ebene-Wellen-Spektrum & das Licht als Summe ebener Wellen: ein abgetastetes
  Winkelspektrum (ein begrenzter Strahl) oder einzelne Wellen bzw.\ Beugungsordnungen. Gitter,
  Schichtstapel und Streukörper wirken auf jede ebene Welle exakt, und zwischen parallelen
  Ebenen ist die freie Ausbreitung exakt (nur eine Phase je Welle). \\
G & Gitterfeld & Feldwerte auf einem Gitter. \\
T & Tabelle & Streumatrix bzw.\ Jones-Tabelle je Ordnung über Wellenlänge und
  Einfallsrichtung. \\
P & radiometrisch & Strahlen ohne Phase (nur inkohärent). \\
\bottomrule
\end{referenz}

Gerechnet werden können derzeit \textbf{R} und – wenn das Python-Paket \code{fieldtracing}
in der Plugin-Umgebung installiert ist – \textbf{H} und \textbf{S}. G, T und P kennt der
Planer bereits; im Transportmenü sind sie ausgegraut (Abschnitt~\ref{sec:workflow-transport}),
und eine Verbindung, die in einem älteren Dokument fest auf eine davon steht, meldet beim Prüfen
bzw.\ Rechnen einen Fehler.

Normalerweise steht der \emph{Transport} einer Verbindung auf \code{auto}: Der \emph{Planer}
wählt für jede Verbindung und Richtung die günstigste Darstellung, die beide Enden annehmen
können, und plant dort, wo ein Block eine andere Darstellung braucht, eine Umwandlung ein
(z.\,B.\ R $\to$ H). \code{auto} wählt nie P, solange eine Darstellung mit Phase möglich ist.
Sie können den Transport auch fest einstellen (Abschnitt~\ref{sec:workflow-transport}), etwa
\code{H} für die Verbindung in eine Kamera nahe einem Fokus oder \code{S} zwischen zwei Gittern.

\paragraph{Wann wird S verwendet?}\index{Spektrum ebener Wellen}
Ein Spektrum ebener Wellen lohnt sich zwischen Blöcken, die ohnehin im Spektrum rechnen. Diese
Blöcke nehmen S \emph{nativ}, also ohne Umwandlung an ihrem Eingang:

\begin{itemize}
  \item Gitter \block{fg.grating} und Schichtstapel \block{fg.stack}: jede Welle für sich durch
        FourierGrating; jede Beugungsordnung geht als eigenes Spektrum weiter,
  \item eine verdichtete Gruppe (\gui{Condensed (plane-wave tables)},
        Abschnitt~\ref{sec:workflow-gruppen}),
  \item der Streukörper \block{fusion.particle}: das Spektrum am Eingang, das rückgestreute
        Fernfeld als Spektrum am Ausgang,
  \item die Detektoren \block{optics.detector} und \block{rtt.detector}: das Bild als Summe der
        Wellen auf den Pixeln,
  \item \block{fusion.convert} mit \param{to}~=~\code{S}.
\end{itemize}

\code{auto} wählt S zwischen zwei Blöcken, die beide S nativ nehmen, sobald das Licht, das den
ersten verlässt, ein Spektrum ist: hinter einem \block{fusion.convert} mit \param{to}~=~\code{S},
hinter einer fest auf \code{S} gestellten Verbindung oder hinter einem Gitter, Stapel, einer
verdichteten Gruppe oder einem Streukörper, die ein Spektrum bekommen haben. Die ganze Kette
nativer Blöcke dahinter bleibt so bei S. Aus Strahlen erzeugt \code{auto} dagegen kein Spektrum
von selbst (R $\to$ S geht über ein Feld und kann sehr viele Wellen ergeben); Workflows ohne
Spektrum behalten daher ihren Plan. S beginnen Sie, indem Sie einen Transport fest auf \code{S}
stellen oder einen \block{fusion.convert} einsetzen, typischerweise
\begin{itemize}
  \item zwischen zwei Gittern oder einem Gitter und einem Stapel: keine Rundreise über ein Feld
        an jedem Block, die Kette ist so genau wie FourierGrating mit beiden Gittern in einer
        Struktur;
  \item zwischen Fokus und Streukörper bzw.\ zwischen Streukörper und Kamera: keine Strahlen
        dazwischen;
  \item in eine Kamera nahe einem Fokus: die Summe der Wellen ist dort das Debye-Integral ohne
        dessen Näherung (weit vom Fokus geht der Detektor selbst über H).
\end{itemize}

Kommt ein Spektrum auf einer Verbindung auf \code{auto} an (etwa aus einem
\block{fusion.convert} oder einem Gitter, das ein Spektrum bekam), bleibt es ein Spektrum, wenn
der nächste Block S nativ nimmt. Vor allen anderen Blöcken wandelt die Ausführung um:
Strahlen werden über ein Feld zum Spektrum (R $\to$ H $\to$ S), ein Feld über die
PFT\footnote{Fourier-Transformation durch stationäre Phase (Fieldtracing).} oder eine FFT (H $\to$
S); ein Spektrum wird vor einem Block, der Felder nimmt, wieder ein Feld (S $\to$ H) und vor einem
Block, der nur Strahlen nimmt, Strahlen (S $\to$ H $\to$ R). Linsen und Spiegel eines Systems
nehmen deshalb nie S. Was ein Spektrum kostet: viele Wellen je Strahl (oft Tausende), bei Gittern
je Welle eine FourierGrating-Lösung -- für breite Bündel hilft \param{field\_transfer}~=~%
\code{interpolated}.

\paragraph{Automatisch Wellenoptik nahe einem Fokus.} Ist im Panel des Workflow-Plugins
\gui{Wave optics near a focus (automatic)} eingeschaltet (Standard), prüft ein Detektor an einer
Verbindung auf \code{auto}, ob sich die Strahlen in seiner Nähe kreuzen oder fokussieren. Dann
rechnet er ab der letzten Fläche mit Feldern, als stünde die Verbindung auf H; sonst bleibt er
bei Strahlen. Eine Verbindung, die fest auf R, H oder S steht, behält ihre Einstellung.

\paragraph{Was wurde verwendet?}\label{sec:workflow-verwendet}
Nach einem Lauf zeigt eine Verbindung auf \code{auto}, was sie tatsächlich getragen hat, wenn es
vom Plan abweicht: \textbf{R$\to$H} heißt geplant R, gerechnet als Feld (Detektor nahe einem
Fokus); \textbf{R$\to$S} hieße geplant R, aber ein Spektrum ist geflossen -- seit \code{auto}
hinter einem Spektrum selbst S plant, kommt das praktisch nicht mehr vor. Der Tooltip nennt in der Zeile \gui{last run} die
Darstellung, die Umwandlungen auf dem Weg und den Grund, z.\,B.\
\code{last run: S (R$\to$H$\to$S) – lens gave rays; spectra on the link; Camera takes them as
they are} -- das auch bei einem festen Transport. \enquote{takes them as they are} heißt: Der
Block rechnet nativ mit dieser Darstellung; sonst steht dort \enquote{converted to rays at the
port of …}. Die Anzeige bleibt, solange das Ergebnis des Blocks am Ende der Verbindung aktuell
ist, und fällt auf die geplante Darstellung zurück, sobald es veraltet.

\paragraph{Hinweise der Spektren.} Die Umwandlungen prüfen ihre Gültigkeit und melden
Abweichungen als Warnung: \code{exec.spectrum\_cut} (ein Teil der Leistung fiel aus dem Spektrum,
meist evaneszente Wellen; die Meldung nennt den Anteil) und \code{exec.spectrum\_rays} (ein
Spektrum wurde zu Strahlen, obwohl das Bündel nicht groß gegen die Wellenlänge ist oder sich
Beugungsordnungen überlappen -- die Strahlen verlieren dann ihre Interferenz). Nach dem Lauf trägt
die betroffene Verbindung bzw.\ der Block ein Warnzeichen (Dreieck mit Ausrufezeichen) neben der
Darstellung, der Tooltip die Meldungen; die Statuszeile nennt die erste, das Plugin-Log alle
(Abschnitt~\ref{sec:workflow-fehler}).

\subsection{Wer rechnet?}

Den Workflow prüft und rechnet das mitgelieferte Plugin \plugin{Workflow}, die
\emph{Workflow-Maschine}\index{Workflow-Maschine}. Sie kennt selbst keinen Rechenkern. Die
Blöcke bringen die Plugins mit, die sie rechnen können:

\begin{itemize}
  \item der \plugin{Raytracer} die Raytatouille-Blöcke \block{rtt.*} (Abschnitte eines
        optischen Systems, Strahlteiler, Spiegel, Detektoren, neue Teile),
  \item \plugin{FourierGrating} das Gitter \block{fg.grating} und den Schichtstapel
        \block{fg.stack},
  \item die Workflow-Maschine selbst die allgemeinen Blöcke: ideale optische Elemente
        (\block{optics.*}), Umwandlungen, Meta- und Anzeigeblöcke.
\end{itemize}

Die Maschine fragt die anderen Plugins nach ihren Blocktypen, nimmt sie in ihre Bibliothek auf
und lässt jeden Block im Prozess seines Plugins rechnen. Fehlt ein Plugin, bleibt sein Block mit
allen Daten im Dokument erhalten; die Prüfung meldet \code{graph.provider\_missing}
(Abschnitt~\ref{sec:workflow-anbieter}).

\subsection{Was wird gerechnet?}\label{sec:workflow-was}

Ähnlich wie bei den Viewer-Knoten in Blender rechnet \gui{Run workflow} nur, was ein
\emph{Anzeigeblock} zeigen soll – und alles, was davor liegt. Kameras und Meta-Blöcke liefern
keine eigenen Ergebnisse: Erst ein \block{display.image} an einer Kamera, ein
\block{display.plot} an einem Meta-Block oder ein \block{display.beam} als Sonde an einem
beliebigen Lichtausgang macht ihr Ergebnis sichtbar. Ohne Anzeigeblock meldet \gui{Run workflow}
\gui{Nothing to show: …} (Abschnitt~\ref{sec:workflow-fehler}).

Ergebnisse früherer Läufe werden aufbewahrt. Ändern Sie einen Parameter, rechnet der nächste
Lauf nur den geänderten Block und alles, was dahinter liegt (Abschnitt~\ref{sec:workflow-status}).

\subsection{Wie entsteht ein Workflow?}\label{sec:workflow-entstehen}

Einen Workflow legt ein Plugin an; ein leeres Dokument hat keinen. Die üblichen Wege:

\begin{description}
  \item[Aus einem optischen System] Laden Sie im Raytracer ein System (\gui{Load system}) und
        wählen Sie \gui{Workflow from system}. Jede Baugruppe (z.\,B.\ ein Objektiv), jeder
        Strahlteiler und jeder Detektor wird ein Block mit allen Daten, und jeder Pfad des
        Systems wird ein Ast. Kapitel~\ref{kap:raytracer}.
  \item[Neu ohne Datei] \gui{New Raytatouille system} (Raytracer) legt eine Quelle mit
        Standardeinstellungen (mm/\si{\micro\metre}, d-Linie, Eintrittspupille 4~mm, ein
        Feldpunkt auf der Achse), eine Blende, einen Detektor und den Ast \code{main} an.
        Dazwischen fügen Sie Teile aus der Bibliothek ein (Abschnitt~\ref{sec:workflow-beispiel}).
\end{description}

Danach bauen Sie den Workflow im Editor weiter aus: Blöcke hinzufügen, ersetzen, gruppieren,
Parameter ändern, Äste aufnehmen. Der Workflow wird mit dem Dokument in der
\datei{.rttf}-Datei gespeichert.

\paragraph{Alte Form in ein Modell umwandeln.} Ein Workflow in der alten Form (ein Graph aus
Sockets, Version~1, wie ihn ältere Versionen oder Skripte schreiben) kennt keine Ports, Äste und
Gruppen: die Liste \gui{Branches} fehlt, und Baugruppen lassen sich nicht öffnen. \gui{Convert to
model} im Workflow-Plugin wandelt ihn um; aus den Lichtwegen des Graphen werden die Äste. Gitter
im Durchlicht, die dort mit \code{out} $\to$ \code{in} hintereinanderhängen, gehen dabei über ihren
hinteren Port \code{back\_out} weiter (RTT-Fusion ab 0.5.1). Die Umwandlung ist ein
Bearbeitungsschritt und lässt sich rückgängig machen.

% ---------------------------------------------------------------------------
\section{Das Fenster \enquote{Node workflow}}\label{sec:workflow-fenster}

Der Editor liegt im Dock \gui{Node workflow} (unten). Es öffnet sich von selbst, sobald ein
Plugin einen Workflow anlegt oder ein Dokument mit Workflow geöffnet wird, und lässt sich mit
\menu{View,Node workflow} ein- und ausblenden. Wie jedes Dock können Sie es abreißen, auf einen
zweiten Bildschirm ziehen oder an einer anderen Seite andocken.

Das Dock ist zweigeteilt (die Trennlinie ist verschiebbar):

\begin{itemize}
  \item links die \emph{Zeichenfläche} mit den Blöcken und Verbindungen,
  \item rechts die Liste \gui{Branches} mit den Ästen und darunter \gui{Steps} mit den
        Schritten des gewählten Asts (Abschnitt~\ref{sec:workflow-aeste}). Bei Workflows in
        der alten Form fehlt dieser Teil.
\end{itemize}

Über der Zeichenfläche steht \gui{Update drawing}: Alle laufenden Plugins, die dem Workflow
folgen, bringen ihre Darstellung sofort auf den Stand des Workflows – der Raytracer zeichnet das
optische System als CAD-Körper –, wie sonst kurz nach jeder Änderung im Editor. Nützlich nach
dem Öffnen eines Dokuments, dessen Workflow noch nicht gezeichnet ist.

Während ein Workflow rechnet, heißt das Dock \gui{Node workflow – running}.

\subsection{So sieht ein Block aus}\label{sec:workflow-blockbild}

Ein Block ist ein Rechteck mit einer farbigen Kopfzeile. Die Farbe zeigt die
\emph{Kategorie}\index{Kategorie} des Blocktyps:

\begin{referenz}{@{}p{0.22\linewidth}p{0.2\linewidth}X@{}}
\toprule
\textbf{Kategorie} & \textbf{Farbe} & \textbf{Beispiele} \\
\midrule
\endhead
\code{source} & grün & Quelle eines Systems, kollimierter Strahl \\
\code{domain} & blau & Linsengruppe, Spiegel, Grenzfläche, Schichtstapel, Gruppe \\
\code{splitter} & violett & Strahlteiler, Gitter \\
\code{detector} & orange & Detektor, Kamera \\
\code{conversion} & grau & Umwandlung zwischen Darstellungen \\
\code{data} & türkis & Datenquellen \\
\code{analysis} & ocker & Kamerabild, Anzeigeblöcke \\
\code{meta}, \code{group} & dunkelgrau & Sweep, Wiederholungen, Optimierung; Gruppen \\
\bottomrule
\end{referenz}

Weitere Kennzeichen:

\begin{itemize}
  \item \textbf{Titel}: der Name des Blocks (ohne Namen seine Kennung), zu lang gekürzt.
  \item \textbf{Ports} (gelb-orange Punkte) für Licht: links die Ports, durch die das Licht beim
        ersten Besuch eines Asts hereinkommt, rechts die übrigen. \textbf{Daten-Anschlüsse}
        (blaue Punkte) stehen links (Eingänge) bzw.\ rechts (Ausgänge). Neben jedem Punkt steht
        sein Name; der Tooltip nennt Art und angenommene Darstellungen bzw.\ Datentyp.
  \item \textbf{Zustand} (Kreis oben links): ein grüner Haken \checkmark{} (\gui{Up to date:
        kept from the last run}) oder ein rotes Kreuz $\times$ (\gui{Out of date: computed again
        at the next run}); siehe Abschnitt~\ref{sec:workflow-status}.
  \item \textbf{Umwandlungen}: Plant der Planer in diesem Block eine Umwandlung, steht oben
        rechts das Zeichen $\rightleftarrows$; der Tooltip listet unter \gui{planned conversions:}
        die Umwandlungen (\code{R $\to$ H: …}).
  \item \textbf{Probleme}: Ein Block mit einem Fehler der letzten Prüfung hat einen
        \emph{roten Rahmen}; der Tooltip nennt Code und Meldung.
  \item \textbf{Läuft}: Der gerade gerechnete Block hat einen hellblauen Rahmen und den Zusatz
        \gui{– running …}.
  \item \textbf{Stumm}: Ein stummgeschalteter Block hat eine graue Kopfzeile und den Zusatz
        \gui{(muted)} (Abschnitt~\ref{sec:workflow-bloecke-bearbeiten}).
\end{itemize}

Der Tooltip eines Blocks zeigt in der ersten Zeile \textit{Kennung}~\code{–}~\textit{Typ} und
darunter seine einfachen Parameter als \code{Name = Wert}.

\subsection{So sehen Verbindungen aus}

\begin{itemize}
  \item Licht-Verbindungen sind gelb-orange, Daten-Links blau. Verläuft eine Verbindung
        rückwärts (das Licht kommt zurück, z.\,B.\ von einem Spiegel), wird sie als Bogen
        unter den Blöcken gezeichnet.
  \item Auf der Mitte einer Verbindung steht für jede Richtung, die ein Ast benutzt, ein
        Schild: ein Pfeil $\blacktriangleright$ bzw.\ $\blacktriangleleft$ (wie auf dem
        Bildschirm; das Schild über dem Band gilt für die zuerst genannte Richtung, das darunter
        für die andere) und die Darstellung – die geplante, z.\,B.\ \code{R}, oder ein fester
        Transport mit Stern, z.\,B.\ \code{H*}. Hat der letzte Lauf etwas anderes verwendet als
        geplant, steht beides da: \code{R$\to$H} oder \code{R$\to$S}
        (Abschnitt~\ref{sec:workflow-verwendet}). Ein Warnzeichen (Dreieck mit Ausrufezeichen)
        dahinter: Die Umwandlung in ein Spektrum oder aus ihm hat etwas gemeldet.
  \item Der Tooltip nennt beide Richtungen mit Blocknamen, z.\,B.\
        \code{Objective A $\to$ Sample: H,} \code{transport fixed H}, und \code{no branch} für eine
        unbenutzte Richtung; nach einem Lauf auch \code{last run: …} mit den Umwandlungen und
        die Hinweise der Spektren.
  \item Eine Verbindung mit einem Problem ist rot gestrichelt.
  \item Fest eingestellte Transporte werden etwas dicker gezeichnet als \code{auto}.
  \item Die Verbindungen des hervorgehobenen Asts sind dicker und hellblau, die Schilder der
        benutzten Richtung nehmen diese Farbe an.
\end{itemize}

Nach einer Änderung bleibt der letzte Plan sichtbar, aber \emph{grau}: Er ist veraltet, bis
\gui{Check workflow} ihn erneuert; der Tooltip trägt den Zusatz \gui{(planned before the last
edit – Check workflow renews it)}.

% ---------------------------------------------------------------------------
\section{Navigation im Editor}\label{sec:workflow-navigation}

Die Zeichenfläche bedienen Sie wie die CAD-Ansichten. Tastenkürzel wirken, solange der Editor
den Fokus hat (einmal hineinklicken), und dann \emph{nicht} auf die CAD-Ansichten.

\begin{referenz}{@{}p{0.34\linewidth}X@{}}
\toprule
\textbf{Aktion} & \textbf{Maus / Tasten} \\
\midrule
\endhead
Verschieben (Pan) & \maus{mittlere Maustaste ziehen}, \maus{rechte Maustaste ziehen},
  \keys{Alt} + \maus{linke Maustaste ziehen}; Trackpad: \maus{mit zwei Fingern wischen} \\
Zoomen am Mauszeiger & \maus{Mausrad}; Trackpad: \maus{Spreizen} oder \keys{Strg} +
  \maus{Wischen mit zwei Fingern} \\
Alles zeigen & \keys{Z} \keys{E} (nacheinander), \keys{Pos1}, \maus{Doppelklick mit der
  mittleren Maustaste}, Trackpad: \maus{Doppeltippen mit zwei Fingern}; Kontextmenü
  \gui{Show all} \\
Auswahl zeigen & \keys{Z} \keys{S}; Kontextmenü \gui{Zoom selected} (ohne Auswahl: alles) \\
Blöcke auswählen & \maus{Klick} auf einen Block; \maus{Rechteck aufziehen} mit der linken
  Maustaste auf leerer Fläche \\
Block verschieben & \maus{Block an der Kopfzeile oder am Rumpf ziehen} \\
Kontextmenü & \maus{Rechtsklick ohne Bewegung} \\
Gruppe verlassen & \keys{Esc} oder Klick auf die Ortsangabe oben links \\
Aufnahme eines Asts abbrechen & \keys{Esc} \\
\bottomrule
\end{referenz}

Ein Rechtsklick, bei dem die Maus bewegt wird, verschiebt nur und öffnet kein Menü. Der Zoom
reicht vom 0,05-Fachen bis zum 8-Fachen.

% ---------------------------------------------------------------------------
\section{Blöcke bearbeiten}\label{sec:workflow-bloecke-bearbeiten}

\subsection{Block hinzufügen}

\begin{schritte}
  \item Rechtsklick auf eine leere Stelle der Zeichenfläche.
  \item \gui{Add block} öffnet ein Untermenü je Kategorie (\code{analysis}, \code{conversion},
        \code{data}, \code{detector}, \code{domain}, \code{meta}, \code{source},
        \code{splitter}); darin stehen die Blocktypen mit ihrer Kennung, der Tooltip
        beschreibt den Typ. Blöcke anderer Plugins erscheinen hier, sobald \gui{Check workflow}
        einmal gelaufen ist. Angeboten werden nur Typen, die diese Installation rechnen kann
        (dazu die Typen, die der Workflow schon enthält); geplante Typen ohne Rechner und
        Blöcke nicht geladener Plugins fehlen in der Liste.
  \item Wählen Sie einen Typ. Der Block erscheint an der Stelle des Rechtsklicks, mit den
        Standardwerten seiner Parameter. Seine Kennung ist der letzte Teil des Typs, z.\,B.\
        \code{image} für \block{display.image}, bei Wiederholung \code{image\_2},
        \code{image\_3} usw.
\end{schritte}

\subsection{Verschieben, Anordnen, Löschen}

Blöcke verschieben Sie mit der Maus; die Position wird mit dem Workflow gespeichert.
\gui{Arrange} (Kontextmenü der Zeichenfläche) ordnet alle Blöcke der aktuellen Ebene neu in
Spalten an: In der neuen Form steht ein Block in der Spalte des ersten Schritts, in dem ein Ast
ihn erreicht. Ein importierter Workflow ohne Positionen wird beim ersten Anzeigen von selbst so
angeordnet.

\gui{Delete block} (Kontextmenü des Blocks) löscht den Block mit seinen Verbindungen und Links.
Äste, die durch ihn liefen, bleiben stehen; die Prüfung meldet sie dann als fehlerhaft, bis Sie
sie löschen oder neu aufnehmen.

\subsection{Stummschalten}

\gui{Mute (bypass)} im Kontextmenü schaltet einen Block stumm, \gui{Unmute} wieder ein. Ein
stummer Block lässt Licht unverändert durch und liefert an Daten-Ausgängen nichts. Bei
Raytatouille-Blöcken werden Pfade bzw.\ Äste, auf denen ein stummer Block liegt, aus der Rechnung
genommen; das Log nennt sie (\code{… left out (muted block on it)}).

\subsection{Block ersetzen}\label{sec:workflow-ersetzen}

\gui{Replace with} tauscht einen Block gegen einen anderen Typ aus und \emph{behält seine
Verbindungen}. Das Untermenü listet nur Typen, bei denen das möglich ist: Jeder verbundene Port
wandert auf den Port gleichen Namens, sonst auf den einzigen Port des neuen Typs. Meta- und
Analyseblöcke stehen nicht zur Wahl. Die Schritte der Äste werden auf die neuen Ports
umgeschrieben; Daten-Links, für die der neue Typ keinen Anschluss hat, entfallen.

Beispiel: Den Probenspiegel eines Linnik-Interferometers ersetzen Sie durch ein Gitter
\block{fg.grating} – der Port \code{front} bleibt verbunden, die Äste laufen nun über das Gitter.
Ersetzt ein Block anderer Art einen Raytatouille-Block, behält er dessen Element und
Durchläufe, sodass der Raytracer ihn weiter zeichnet.

\subsection{Aktionen der Plugins an Blöcken}

Manche Plugins bieten Befehle für bestimmte Blocktypen an. Sie stehen im Kontextmenü des Blocks
unter \gui{Parameters…}, z.\,B.:

\begin{itemize}
  \item \gui{Section $\to$ group} (Raytracer) an einem \block{rtt.section}: macht aus einem
        Abschnitt eines Systems eine Gruppe mit einem Block je Element
        (Abschnitt~\ref{sec:workflow-gruppen}).
  \item \gui{Store panel configuration in block} und \gui{Load block configuration into panel}
        (FourierGrating) an \block{fg.grating} und \block{fg.stack}
        (Kapitel~\ref{kap:fouriergrating}).
  \item \gui{Store job in block} und \gui{Load block job into panel} (hp-FEM) an
        \block{hpfem.grating} (Kapitel~\ref{kap:hpfem}).
\end{itemize}

Diese Befehle laufen als normale Plugin-Jobs (Abschnitt~\ref{sec:plugin-jobs}) und sind ein
Rückgängig-Schritt.

% ---------------------------------------------------------------------------
\section{Verbinden}\label{sec:workflow-verbinden}

\subsection{Zwei Ports verbinden}

Ziehen Sie mit der linken Maustaste von einem Port zu einem Port eines anderen Blocks. Während
des Ziehens folgt eine gestrichelte Linie der Maus. Lassen Sie über dem Ziel-Port los, entsteht
die Verbindung. Regeln:

\begin{itemize}
  \item Ein Port hat höchstens \emph{eine} Verbindung. Ist einer der beiden Ports schon belegt,
        passiert nichts. Verzweigen und Zusammenführen geschieht in Blöcken (Strahlteiler).
  \item Ein Port lässt sich nicht mit einem Daten-Anschluss verbinden.
  \item Ziehen Sie einen Port auf einen \emph{Anzeigeblock} (z.\,B.\ \block{display.beam}),
        entsteht keine Verbindung, sondern ein \emph{Beobachter-Link}: Die Anzeige zeigt das Licht,
        das diesen Port verlässt; ein Ausgang darf mehrere Links speisen.
\end{itemize}

\subsection{Daten verbinden}

Daten-Anschlüsse bleiben gerichtet: Ziehen Sie von einem Ausgang zu einem Eingang (oder
umgekehrt). Beide müssen von derselben Art sein; zwei Eingänge oder zwei Ausgänge lassen sich nicht
verbinden. Ein Eingang nimmt nur \emph{einen} Link an. Typische Daten-Links:

\begin{itemize}
  \item \code{image} eines Detektors $\to$ \code{in} eines \block{display.image},
  \item \code{results} eines Meta-Blocks $\to$ \code{in} eines \block{display.plot},
  \item \code{drive} eines Meta-Blocks $\to$ \code{drive} eines Spiegels, Abschnitts oder
        Strahlteilers (Verschiebung längs der Flächennormalen, wie ein Piezo-Phasenschieber)
        oder $\to$ \code{pose} eines Blocks (Lageänderung),
  \item der zu beobachtende Ausgang (z.\,B.\ \code{camera.image}) $\to$ \code{observe} des
        Meta-Blocks.
\end{itemize}

Jeder Block mit Lage – Quellen, optische Elemente, Strahlteiler, Detektoren – hat dazu den
Daten-Eingang \code{pose} (sechs Werte: Verschiebung \code{dx}, \code{dy}, \code{dz} in mm und
Drehung \code{rx}, \code{ry}, \code{rz} in rad).

\subsection{Verbindung löschen}

Rechtsklick auf die Verbindung $\to$ \gui{Delete connection}, bei einem Daten-Link
\gui{Delete link}. Verbindungen lassen sich auch mit der Maus anklicken (sie werden weiß
hervorgehoben); die Trefferzone ist breiter als die Linie.

\subsection{Transport festlegen}\label{sec:workflow-transport}\index{Transport}

Rechtsklick auf eine Verbindung zweier Ports bietet drei Untermenüs:

\begin{itemize}
  \item \gui{Transport (both directions)} – für beide Richtungen,
  \item \gui{Transport} \textit{A}~$\to$~\textit{B} – nur für die Richtung von Block A nach B,
  \item \gui{Transport} \textit{B}~$\to$~\textit{A} – nur für die Gegenrichtung.
\end{itemize}

Jedes bietet \gui{auto – the planner chooses}, \gui{R – rays}, \gui{H – fields (hybrid,
Fieldtracing)}, \gui{S – plane-wave spectra}, \gui{G – gridded fields}, \gui{T – tables (Jones,
S-matrix)} und \gui{P – radiometric (no phase)}; der aktuelle Transport ist abgehakt. Bei einem
Link der alten Form heißt das Untermenü einfach \gui{Transport}. \code{auto} überlässt die Wahl dem
Planer; jeder andere Buchstabe legt die Darstellung fest (Schild mit Stern, z.\,B.\ \code{S*}).
Ausgegraut ist, was Ihre Installation nicht rechnen kann; der Tooltip sagt warum: H und S
brauchen \code{fieldtracing}, G, T und P werden auf Verbindungen noch nicht gerechnet. Welche
Darstellungen verfügbar sind, trägt \gui{Check workflow} ins Dokument ein.

\begin{tipp}
Zwei Gitter (oder Gitter und Stapel) hintereinander: Stellen Sie die Verbindung zwischen ihnen
auf \code{S}. Das Spektrum geht ohne Umweg über ein Feld von Gitter zu Gitter; vor der nächsten
Linse wandelt die Ausführung von selbst in Strahlen um. Ob und wie umgewandelt wurde, zeigt nach
dem Lauf der Tooltip der Verbindungen (\gui{last run}).
\end{tipp}

\begin{tipp}
Für ein beugungsbegrenztes Bild nahe einem Fokus (Pupillen- oder Fokuskamera) stellen Sie die
Verbindung in die Kamera auf \code{H} und wählen im Workflow-Panel ein kleines
\gui{Image window (half width, mm)}, z.\,B.\ 0.05. Meist genügt es aber, die Verbindung auf
\code{auto} zu lassen und \gui{Wave optics near a focus (automatic)} eingeschaltet zu haben.
\end{tipp}

% ---------------------------------------------------------------------------
\section{Äste}\label{sec:workflow-aeste}\index{Ast}

Ein \emph{Ast} ist ein benannter Lichtweg: von einer Quelle durch eine Folge von Blöcken bis zu
einem Detektor. Jeder Schritt nennt den Block, den Port, durch den das Licht hereinkommt, und den
Port, durch den es ihn verlässt. Ein Spiegel wird durch denselben Port verlassen, durch den das
Licht kam; ein Strahlteiler durch den Port der gewünschten Richtung. \emph{Nur die Äste werden
gerechnet} – Verbindungen allein bestimmen noch keinen Lichtweg. \gui{Workflow from system} legt
je Pfad des optischen Systems einen Ast an.

\subsection{Die Liste \enquote{Branches}}

Rechts neben der Zeichenfläche:

\begin{description}
  \item[\gui{Branches}] alle Äste der aktuellen Ebene (Tooltip: \gui{Exactly these branches are
        computed. Select one to highlight it.}). Ein Klick hebt den Ast hervor: Seine Blöcke
        bekommen einen hellblauen Rahmen, seine Verbindungen werden hellblau.
  \item[\gui{Steps}] die Schritte des gewählten Asts, je Zeile
        \textit{Block}\code{:}~\textit{Port~ein}~$\to$~\textit{Port~aus}, ggf.\ mit dem
        Kanal in eckigen Klammern; $\cdot$ steht für \enquote{keiner} (die Quelle hat keinen
        Eingang, der Detektor keinen Ausgang).
  \item[\gui{New branch}] beginnt die Aufnahme eines Asts (während der Aufnahme heißt der Knopf
        \gui{Cancel}).
  \item[\gui{Delete}] löscht den markierten Ast.
\end{description}

Innerhalb einer Gruppe steht über der Liste \gui{Inside:} mit dem Namen der Gruppe, und die
Liste zeigt die \emph{inneren} Äste der Gruppe.

\subsection{Einen Ast aufnehmen}\label{sec:workflow-ast-aufnehmen}

\begin{schritte}
  \item Klicken Sie auf \gui{New branch} (oder im Kontextmenü der Zeichenfläche auf
        \gui{New branch…}). Unter der Liste erscheint der Hinweis \gui{New branch: click the
        port of a source the light leaves by (Esc cancels)}.
  \item Klicken Sie den Port der \emph{Quelle}, durch den das Licht sie verlässt.
  \item Klicken Sie nacheinander bei jedem folgenden Block den Port, durch den das Licht ihn
        \emph{verlässt} (Hinweis: \gui{Click the port the light leaves the next block by, or the
        detector (Esc cancels)}). Der nächste Block ist immer der am anderen Ende der Verbindung
        des zuletzt gewählten Ports; Klicks auf andere Blöcke werden ignoriert. Bei einem Spiegel
        klicken Sie also den Port, durch den das Licht hereinkam.
  \item Klicken Sie zuletzt den \emph{Detektor} (den Block oder seinen Port). Der Hinweis lautet
        jetzt \gui{The branch reaches a detector}.
  \item Der Dialog \gui{New branch} fragt nach dem \gui{Name of the branch:} (Vorschlag
        \code{branch 1}, \code{branch 2}, …). Der Name muss eindeutig sein.
\end{schritte}

Der aufgenommene Teil wird während der Aufnahme gelb hervorgehoben. \keys{Esc} oder
\gui{Cancel} brechen ab. Führt der Ast durch eine Gruppe, wird automatisch der innere Ast der
Gruppe eingetragen, der zu den beiden Ports passt.

\subsection{Kanäle eines Schritts}\label{sec:workflow-kanaele}

Ein Schritt kann einen \emph{Kanal} tragen, der genauer sagt, wie das Licht den Block durchläuft.
Doppelklicken Sie in \gui{Steps} auf den Schritt; der Dialog \gui{Channel of the step} bietet
Vorschläge an, Sie können aber auch frei tippen:

\begin{itemize}
  \item \gui{(none)} – kein Kanal,
  \item bei einer Gruppe: die inneren Äste, die zu den Ports des Schritts passen,
  \item bei einem Raytatouille-Block: \code{ghost:}\textit{Fläche} – ein \emph{Geisterbild},
        das an dieser Fläche reflektiert wird (eine ungerade Zahl von Flächen verlässt den Block
        durch den Port, durch den das Licht kam, eine gerade durch den anderen),
  \item beliebiger Text, z.\,B.\ eine Beugungsordnung.
\end{itemize}

% ---------------------------------------------------------------------------
\section{Gruppen und Varianten}\label{sec:workflow-gruppen}\index{Gruppe}\index{Variante}

Eine \emph{Gruppe} fasst mehrere Blöcke zu einem zusammen – etwa die Linsen eines Objektivs. Nach
außen hat sie die Ports, über die Verbindungen in sie hineinführen; innen hat sie eigene Äste.
Eine Gruppe kann mehrere \emph{Varianten} haben, zwischen denen Sie umschalten, z.\,B.\ ein
reales Objektiv und eine ideale Linse.

\subsection{Gruppe bilden}

\begin{schritte}
  \item Wählen Sie die Blöcke aus (Rechteck aufziehen oder anklicken).
  \item Rechtsklick auf einen der Blöcke $\to$ \gui{Group selected blocks…}.
  \item Geben Sie im Dialog \gui{Group} den \gui{Name of the group:} ein (Vorschlag
        \code{Group}).
\end{schritte}

Verbindungen über den Rand der Gruppe werden zu exportierten Ports
(\textit{Block}\code{.}\textit{Port}), jedes Stück eines Asts innerhalb der Gruppe wird ein
innerer Ast, und der Schritt durch die Gruppe nennt ihn als Kanal. Abgelehnt wird – mit
Begründung in einem Meldungsfenster –, wenn

\begin{itemize}
  \item die Auswahl leer ist oder der Workflow die alte Form hat
        (\gui{Select blocks of a model first.}),
  \item eine Quelle oder ein Detektor dabei ist (\gui{… is a source or a detector: branches
        begin and end outside a group.}),
  \item ein Daten-Link den Rand kreuzt (\gui{The data link … crosses the group's border; group
        the blocks it joins too, or remove it.}),
  \item ein Ast die Gruppe nicht über Verbindungen betritt und verlässt.
\end{itemize}

\subsection{In eine Gruppe hineingehen}

Ein Doppelklick auf eine Gruppe (oder \gui{Open group}) zeigt ihr Inneres. Oben links steht die
Ortsangabe, z.\,B.\ \code{Workflow $\blacktriangleright$ Objective (real)  – Esc: back}, mit
dem Namen der Gruppe und der aktiven Variante in Klammern. \keys{Esc}, ein Klick auf die
Ortsangabe oder \gui{Back out of the group} (Kontextmenü der Zeichenfläche) führen eine Ebene
zurück. Gruppen können geschachtelt sein.

Ein außen hervorgehobener Ast setzt sich innen fort: Beim Hineingehen wird der innere Ast
hervorgehoben, den sein Schritt durch die Gruppe nennt; beim Verlassen wieder der äußere.

\subsection{Varianten}

Das Untermenü \gui{Variant} im Kontextmenü einer Gruppe listet die Varianten; ein Häkchen
markiert die aktive. Ein Klick macht eine andere Variante aktiv.

\begin{description}
  \item[\gui{New variant…}] kopiert die aktive Variante unter einem neuen Namen (Dialog
        \gui{New variant}, \gui{A copy of "…", named:}, Vorschlag \textit{Name}~\code{copy})
        und macht die Kopie aktiv. Darin ersetzen Sie dann z.\,B.\ die Linsengruppe durch eine
        ideale Linse.
  \item[\gui{Rename variant…}] benennt die aktive Variante um; ein leerer oder vergebener Name
        wird mit \gui{The name is empty or taken.} abgelehnt.
  \item[\gui{Delete variant "…"}] löscht die aktive Variante; die letzte bleibt immer.
\end{description}

Alle Varianten müssen dieselben Ports und Parameter nach außen geben, sonst meldet die Prüfung
\code{graph.group\_ports} bzw.\ \code{graph.group\_params}.

\subsection{Verdichten und Auflösen}

\gui{Condensed (plane-wave tables)} (Häkchen) rechnet eine Gruppe nicht Block für Block, sondern
als einen Block, der je innerem Ast eine Tabelle über einfallende ebene Wellen anwendet. Das
lohnt sich für seitlich unveränderliche Gruppen wie Schichtstapel oder Platten; eine Gruppe,
deren Wirkung vom Auftreffort abhängt (eine Linse) oder die s- und p-Polarisation mischt, lässt
sich so nicht verdichten (Fehler \code{exec.condense} beim Rechnen).

\gui{Ungroup} löst die Gruppe auf und setzt ihre Blöcke wieder auf die äußere Ebene.

\subsection{Abschnitt in Gruppe verwandeln}

Ein importierter Raytatouille-Abschnitt (\block{rtt.section}) enthält oft mehrere Elemente einer
Baugruppe, z.\,B.\ alle Linsen eines Objektivs. \gui{Section $\to$ group} im Kontextmenü (vom
Raytracer) macht daraus eine Gruppe mit einem Block je Element (Linsengruppe, Blende …), jeder
mit seinen Daten und seinem Teil jedes Durchlaufs; die Durchläufe des Abschnitts werden die inneren
Äste. Das Licht bleibt dasselbe. Der Befehl ist ein Rückgängig-Schritt.

% ---------------------------------------------------------------------------
\section{Parameter der Blöcke}\label{sec:workflow-parameter}

Ein Doppelklick auf einen Block (außer Gruppen und Anzeigeblöcken) oder \gui{Parameters…} im
Kontextmenü öffnet den Dialog \gui{Parameters of} \textit{Name}. Er ist in \emph{Reiter}
gegliedert: zuerst \gui{Overview}, ein Übersichtsbild des Blocks
(Abschnitt~\ref{sec:workflow-uebersicht}), dann die Parameter des Blocktyps
(Anhang~\ref{kap:workflow-bloecke}) und – bei Blöcken mit eigenen Daten – die Felder, die das
zuständige Plugin für diesen Block beschreibt.

\begin{itemize}
  \item Der Reiter \gui{Block} enthält die Parameter des Typs. Hat ein Block Datenfelder
        (z.\,B.\ ein importierter Raytatouille-Block), folgt ein Reiter
        je Element bzw.\ Gruppe von Daten – Radien, Konstanten der Kegelschnitte, Aperturen,
        Lagen der Flächen, Gläser, Wechselwirkungen. Längen stehen in mm, Winkel in Grad; eine
        Größe, die in der Optimierung \enquote{variabel} ist, erscheint als Kontrollkästchen.
  \item Die Quelle eines Systems zeigt die Systemeinstellungen (Wellenlängen, Felder, Apertur,
        Objekt).
  \item Die Einheit steht in Klammern hinter der Beschriftung, z.\,B.\ \gui{focal\_length (mm)}.
  \item Zahlen tippen Sie mit Dezimalpunkt; Listen und Punkte als durch Kommas getrennte
        Zahlen (\gui{comma-separated numbers}), z.\,B.\ \code{0, 0, 1}; Dateien mit \gui{…}.
  \item Der Tooltip eines Felds beschreibt den Parameter.
\end{itemize}

Mit \gui{OK} werden \emph{nur die Werte geschrieben, die Sie geändert haben}. Ein Wert mit mehr
Stellen, als der Dialog zeigt, bleibt dadurch exakt erhalten, solange Sie ihn nicht anfassen. Hat
ein Block keine Parameter, steht dort \gui{This block has no parameters.}

Interne Datenparameter (z.\,B.\ das gespeicherte System einer Quelle) zeigt der Dialog nicht.

\subsection{Die Übersicht}\label{sec:workflow-uebersicht}\index{Uebersichtsbild@Übersichtsbild}\index{Overview}

Der erste Reiter, \gui{Overview}, zeigt das Teilmodell des Blocks als Zeichnung: maßstäblich, mit
Achsen in mm bzw.\ µm, Glas- und Materialfarben samt Legende, Bemaßungen und Hinweisen unter dem
Bild. Gezeichnet wird mit den Werten, die \emph{im Dialog} stehen – auch solchen, die Sie noch
nicht mit \gui{OK} übernommen haben. So sehen Sie, was eine Änderung an Radius, Dicke oder
Schichtfolge bewirkt, bevor Sie sie übernehmen.

\begin{description}
  \item[Raytatouille-Blöcke] (\block{rtt.section}, \block{rtt.mirror}, \block{rtt.splitter},
        \block{rtt.detector}, \block{rtt.part.*}): der Meridionalschnitt (y–z-Ebene, $z$ nach
        rechts) der Elemente des Blocks. Jede Linse erscheint als geschlossene Kontur ihrer
        Glasteile aus den exakten Pfeilhöhen der Flächen über ihren Aperturen, gefärbt nach Glas;
        die Legende nennt das Glas und seine Brechzahl bei der Referenzwellenlänge. Kittglieder
        zeigen je Glas eine Fläche. Spiegel, ideale Linsen (mit Pfeilspitzen und $f$), Strahlteiler
        (gestrichelt) und Detektoren (blau) erscheinen als Profil, eine Blende als Marken neben
        ihrer Öffnung mit ihrem Durchmesser, die optische Achse strichpunktiert. Bei einem
        Abschnitt läuft ein Fächer paralleler Strahlen den ersten Lauf des Blocks entlang. Bei
        einer einzelnen Linse nennen die Hinweise Radien, Dicke und Durchmesser; bemaßt ist die
        Mittendicke bzw.\ die Länge vom ersten bis zum letzten Scheitel. Die Koordinaten sind die
        der Baugruppe, in der der Block liegt.
  \item[Quelle] (\block{rtt.source}): das ganze System im Meridionalschnitt, ohne Strahlen, mit
        Wellenlängen und Zahl der Feldpunkte.
  \item[FourierGrating-Blöcke] (\block{fg.grating}, \block{fg.stack}): der Querschnitt der
        Konfiguration in der x–y-Ebene (µm) über zwei Perioden: Substrat unten, die Schichten von
        unten nach oben, Superstrat oben, jede Fläche in der Farbe ihres Materials (Legende mit
        $n$ bzw.\ $n + \mathrm{i}k$ oder dem Namen des Modells). Rechteck-, Sinus-, Dreiecks- und
        Stützstellenprofile werden exakt gezeichnet, 2D-Schichten (\code{polygon2d},
        \code{contour2d}) im Schnitt auf Höhe ihres ersten Einschlusses, unbekannte Schichttypen
        (\code{cad}, \code{json}) schraffiert mit ihrem Namen. Rechts sind die Schichtdicken
        bemaßt, oben die Periode $\Lambda$; ein Pfeil zeigt den Einfall ($\lambda$, $\theta$, ggf.\
        $\psi$, TE/TM), bei einem Gitter gestrichelt die nullte reflektierte bzw.\ transmittierte
        Ordnung. Sind Dicke und Breite sehr verschieden, werden die Achsen verschieden skaliert
        (Hinweis unter dem Bild). Ein Gitter, das nur eine Jones-Tabelle hält, erscheint als
        Probenebene mit seinen Ordnungen.
  \item[Ideale Optik] (\block{optics.*}, gezeichnet vom Workflow-Plugin): Linsen als Symbol mit
        Brennpunkten und $f$, Detektoren als Ebene mit Größe und ankommendem Licht, die
        kollimierte Quelle als Strahlenbündel, Spiegel mit schraffierter Rückseite, Strahlteiler
        gestrichelt, Grenzflächen mit beiden Brechzahlen; das Teilchen
        (\block{fusion.particle}) als Kreis in µm.
  \item[Alle anderen Blöcke] (z.\,B.\ Anzeige-, Meta- und Umwandlungsblöcke) zeigen ihr
        \emph{Schema}: den Block mit seinen Lichtports (links hinein, rechts hinaus),
        Datenanschlüssen, Typ, Beschreibung und den ersten Parameterwerten.
\end{description}

\paragraph{Aktualisieren.} Das Bild wird gezeichnet, sobald der Reiter sichtbar ist, und
erneut etwa 0,7~Sekunden nach der letzten Änderung eines Werts im Dialog; \gui{Update} zeichnet
sofort. Während das Plugin rechnet, steht oben rechts \gui{computing the overview…}; das
vorige Bild bleibt so lange stehen. Das Plugin wird dafür gestartet, falls es noch nicht läuft
(beim ersten Mal kann das einige Sekunden dauern). Der Dialog merkt sich für die Sitzung, ob Sie
ihn zuletzt auf der Übersicht oder auf einem Parameterreiter geschlossen haben, und öffnet den
nächsten ebenso; auf einem Parameterreiter wird erst gezeichnet, wenn Sie die Übersicht wählen.

\paragraph{Bild kopieren oder speichern.} Das Kontextmenü des Bilds bietet \gui{Copy image}
(in die Zwischenablage) und \gui{Save image…} (PNG in doppelter Auflösung).

\paragraph{Wenn kein Bild kommt.} Kann ein Plugin den Block nicht zeichnen, steht sein Grund in
Rot über dem Schema, z.\,B.\ \gui{the Raytracer's overview needs the Python package raytatouille}
oder \gui{… holds no FourierGrating configuration (Store panel configuration in block, or name a
configuration file)}. Fehlen die Glaskataloge eines Systems, wird die Geometrie trotzdem
gezeichnet; der Hinweis \gui{glasses unknown (…) – geometry only, no rays} nennt den Grund
(Glaskataloge: Abschnitt~\ref{sec:workflow-eingebettet}). Das Bild ist mit dem Dialog verbunden;
schließen Sie ihn, bevor die Antwort kommt, wird sie verworfen.

\begin{hinweis}
Nach jeder Änderung im Editor ruft die Anwendung nach etwa 0,6~Sekunden die Plugins auf, die dem
Workflow folgen, sofern sie laufen und nichts zu tun haben. Der Raytracer zeichnet so mit
\gui{Follow the workflow} sein System im CAD-Modell neu, wenn sich Radien, Lagen usw.\ geändert
haben. Ein halb eingegebener Wert wartet einfach auf die nächste Änderung.
\end{hinweis}

\begin{hinweis}
Jede Änderung im Node-Editor selbst (Blöcke, Verbindungen, Parameter, Äste, Gruppen) ist ein
Rückgängig-Schritt (\enquote{Edit node workflow}); das Verschieben eines Blocks zählt einmal,
beim Loslassen. Auch die Läufe von Plugins, die den Workflow geändert haben (z.\,B.\
\gui{Workflow from system}, \gui{Section $\to$ group}, eine Optimierung), lassen sich rückgängig
machen: \menu{Edit,Undo} setzt den Workflow auf den Stand vor diesem Lauf zurück.
\end{hinweis}

% ---------------------------------------------------------------------------
\section{Das Workflow-Plugin}\label{sec:workflow-plugin}

Geprüft und gerechnet wird mit dem Plugin \plugin{Workflow} (Plugin-Dock, Plugins-Menü oder
Werkzeugspalte). Es hat zwei Aktionen und ein Panel mit den Bildeinstellungen.

\subsection{Die Einstellungen im Panel}\label{sec:workflow-panel}

\begin{referenz}{@{}>{\raggedright\arraybackslash}p{0.34\linewidth}>{\raggedright\arraybackslash}p{0.17\linewidth}X@{}}
\toprule
\textbf{Feld} & \textbf{Standard} & \textbf{Bedeutung} \\
\midrule
\endhead
\gui{Rays per field and wavelength} & 20000 (7 … 500000) & Zahl der Strahlen je Feldpunkt
  und Wellenlänge für die Quellen des Workflows (kohärentes Bild). Die Strahlen liegen auf einem
  hexapolaren Raster; gewählt wird die Zahl der Ringe $r$, deren $1 + 3r(r+1)$ Strahlen dem Wert
  am nächsten kommen. Der Speicherbedarf wächst mit Strahlen $\times$ Ereignissen, weil jeder
  Strahlweg aufgezeichnet wird. \\
\gui{Pixels (long side)} & 512 (16 … 4096) & Pixelzahl der Bilder an der langen Seite. \\
\gui{Image window (half width, mm)} & 0 & Halbe Breite des Bildfensters aller
  Raytatouille-Detektoren; 0: jeder Detektor mit seiner eigenen Größe. Für einen Fokus, der als
  Feld gerechnet wird, ein kleines Fenster, z.\,B.\ 0.05~mm. \\
\gui{Gaussian illumination (w / pupil radius)} & 0 (0 … 10) & 0: gleichmäßige Ausleuchtung
  (ein hart begrenztes Bündel, dessen Fresnel-Beugung eine Feld-Verbindung zeigt); sonst
  Gauß-Profil mit diesem Verhältnis von Taillenradius zu Pupillenradius, z.\,B.\ 0.5 für einen
  Laserstrahl. \\
\gui{Wave optics near a focus (automatic)} & an & Detektoren an Verbindungen auf \code{auto}
  rechnen nahe einem Fokus oder einer Kaustik mit Feldern (Abschnitt~\ref{sec:workflow-darstellungen}).
  Braucht \code{fieldtracing}. \\
\gui{Field accuracy $\varepsilon$ (H links)} & 0.02 (0.001 … 0.2) & Erlaubter Fehler einer
  Feldausbreitung (Verbindung auf H), bezogen auf die Spitzenamplitude. Fieldtracing wählt das
  günstigste Verfahren innerhalb dieser Grenze; nahe einem Fokus erlaubt 0.05 oft
  Debye (Sekunden) statt Rayleigh-Sommerfeld (Minuten bis Stunden). Das Log nennt Verfahren und
  Fehler. \\
\bottomrule
\end{referenz}

Diese Einstellungen gehören zum Dokument. Ändern Sie eine davon, gelten alle Ergebnisse, die von
ihr abhängen, als veraltet.

\subsection{Check workflow}\label{sec:workflow-pruefen}

\gui{Check workflow} (\enquote{Check the node workflow and mark problems; the plan per link})
prüft den Workflow, ohne zu rechnen:

\begin{enumerate}
  \item Es fragt die anderen Plugins nach ihren Blocktypen und ergänzt die Bibliothek
        (\gui{Add block}).
  \item Es übernimmt die Glaskataloge aus dem Raytracer-Panel und speichert Dateien, die Blöcke
        nennen, in den Blöcken (Abschnitt~\ref{sec:workflow-eingebettet}).
  \item Es prüft Blöcke, Ports, Verbindungen, Äste, Gruppen und Varianten. Für jeden (nicht
        stummgeschalteten) Block, den diese Installation nicht rechnen kann, gibt es die
        Warnung \code{check.no\_evaluator}.
  \item Ist alles fehlerfrei, wird der Workflow entlang der Äste \enquote{abgerollt}, und der
        Planer wählt für jede Verbindung und Richtung die Darstellung.
  \item Es bestimmt, welche Blöcke noch aktuelle Ergebnisse haben (Haken und Kreuze).
\end{enumerate}

Danach sind Probleme im Editor markiert (rote Rahmen, rot gestrichelte Verbindungen, Tooltips),
und die Verbindungen zeigen den neuen Plan. Das Plugin-Log listet jede Meldung als
\textit{Schwere}\code{:}~\textit{Code}~\code{–}~\textit{Meldung}~\code{(}\textit{Ort}\code{)},
die Statuszeile fasst zusammen:

\begin{lstlisting}
Workflow: workflow: 0 error(s), 1 warning(s); planned RH, 2 conversion(s)
\end{lstlisting}

Hat das Dokument keinen Workflow, meldet die Aktion \gui{No node workflow yet – a plugin makes
one (Raytracer: Workflow from system, New Raytatouille system).}

\subsection{Run workflow}\label{sec:workflow-ausfuehren}

\gui{Run workflow} (\enquote{Compute what the display blocks of the node workflow show: camera
images, sweep curves, beam probes; links set to H as fields}) rechnet:

\begin{enumerate}
  \item Wie beim Prüfen werden Blockkataloge geholt und Dateien eingebettet. Hat der Workflow
        Fehler, bricht der Lauf ab: \gui{The workflow has errors – Check workflow marks them:
        …}.
  \item Alle Anzeigeblöcke werden der Reihe nach gerechnet, jeweils mit allem, was davor liegt;
        gemeinsame Blöcke nur einmal. Der Fortschritt steht in der Statuszeile
        (\code{display:}~\textit{Name}, dann je Pfad, Feldpunkt, Wellenlänge und Wert), der
        gerade gerechnete Block hat einen hellblauen Rahmen.
  \item Eine Anzeige, deren gespeichertes Ergebnis zum aktuellen Stand passt, wird
        übersprungen (\code{(unchanged)}).
  \item Jede Anzeige wird ein Ergebnis im neuen Ergebnissatz
        \code{Workflow}~\textit{Systemname} (bzw.\ \code{Workflow run}) – sichtbar im Dock
        \gui{Results} und im Fenster der Anzeige (Abschnitt~\ref{sec:workflow-anzeigen}).
  \item Die Blöcke bekommen ihren Zustand, die Verbindungen ggf.\ die tatsächlich verwendete
        Darstellung (\code{R$\to$H}).
\end{enumerate}

Die Statuszeile meldet am Ende, was gerechnet wurde, z.\,B.\
\code{Workflow: workflow: Pupil camera, Image (unchanged)}. Wurde nichts gerechnet, steht dort
\code{workflow: nothing computed –} und der erste Grund (\code{(and n more in the log)}).
Warnungen, etwa zu wenige Pixel für einen kleinen Fokus, hängen als \code{– warning: …} an.

\paragraph{Abbrechen.} \gui{Cancel} oder \keys{Esc} wirken zwischen den Rechenschritten
(Pfad, Feldpunkt, Wellenlänge, Wert). Bereits fertige Anzeigen behalten ihr Ergebnis.

\begin{tipp}
Ein Kamerabild bleibt schwarz, wenn das Licht in weniger als ein Pixel fokussiert. Das Log
schlägt dann vor, \gui{Wave optics near a focus} einzuschalten oder die Verbindung in die Kamera
auf H zu stellen und ein kleines Bildfenster zu wählen. Deckt ein Bild zu wenige Pixelmitten ab,
nennt die Warnung die Größe des Flecks und ein passendes \gui{Image window} – oder empfiehlt, die
\gui{Pixels} zu erhöhen.
\end{tipp}

% ---------------------------------------------------------------------------
\section{Statusanzeigen und Cache}\label{sec:workflow-status}\index{Cache}

Die Workflow-Maschine bewahrt die Zwischenergebnisse früherer Läufe auf (bis etwa 1~GB; die
ältesten fallen zuerst heraus, mit einer Notiz im Plugin-Log). Jedes Ergebnis ist unter einem
\emph{Fingerabdruck} abgelegt: einer Prüfsumme aus Blocktyp, Parametern, Transporten,
Einstellungen des Panels, Glaskatalogen und den Fingerabdrücken von allem, was davor liegt. Ein
gleicher Fingerabdruck bedeutet ein gleiches Ergebnis.

\begin{itemize}
  \item \textbf{Grüner Haken} \checkmark{} an einem Block: Sein Ergebnis ist aktuell und wird
        beim nächsten Lauf nicht neu gerechnet.
  \item \textbf{Rotes Kreuz} $\times$: Das Ergebnis ist veraltet und wird beim nächsten Lauf
        neu gerechnet.
  \item Kein Kreis: Für diesen Block gibt es (noch) keinen Zustand, z.\,B.\ vor dem ersten
        Prüfen.
\end{itemize}

Eine Änderung im Editor färbt den geänderten Block und \emph{alles dahinter} sofort rot: entlang
der Äste, über Daten-Links und die inneren Abhängigkeiten der Blöcke. Eine Ausnahme ist der
\code{drive}-Ausgang eines Meta-Blocks: Er gibt zwischen den Läufen nichts weiter. Der nächste
Lauf rechnet nur die roten Blöcke; \gui{Check workflow} und jeder Lauf melden den Zustand neu.

\begin{hinweis}
Der Cache liegt im Prozess des Workflow-Plugins. Wird das Plugin neu gestartet
(\menu{Plugins,Reload plugins}, neuer Interpreter), ist er leer. Die Ergebnisse der Anzeigen
selbst sind aber im Dokument gespeichert (Abschnitt~\ref{sec:workflow-gespeichert}); eine
unveränderte Anzeige gilt daher weiter als aktuell.
\end{hinweis}

% ---------------------------------------------------------------------------
\section{Anzeigefenster}\label{sec:workflow-anzeigen}\index{Anzeigeblock}

Jeder Anzeigeblock hat ein eigenes Fenster. Ein \emph{Doppelklick} auf den Block oder
\gui{Show window} im Kontextmenü öffnet es; die Parameter des Anzeigeblocks erreichen Sie dann
über \gui{Parameters…} im Kontextmenü. Das Fenster öffnet sich schwebend mit dem Titel
\gui{Display:}~\textit{Name}; ziehen Sie seine Titelleiste an eine Seite des Hauptfensters, um es
anzudocken. Es zeigt immer das neueste Ergebnis dieser Anzeige.

\begin{description}
  \item[\block{display.image}] zeigt ein Bild – das Ergebnis eines Detektors oder einer
        Kamera – wie das Image-Fenster
        (Abschnitt~\ref{sec:image-fenster}): Achsen in mm, Farbbalken, \gui{log}, Ablesen,
        Zoomen, Verschieben.
  \item[\block{display.plot}] zeigt Kurven wie das Plot-Dock (Abschnitt~\ref{sec:plot-dock}):
        an einem Meta-Block die Messgröße über dem variierten Parameter, an einem Bild seine
        mittlere Zeile und Spalte, an Zahlen ihre Werte über dem Index.
  \item[\block{display.beam}] ist eine \emph{Sonde}: Es zeigt das Licht an einem beliebigen
        Ausgang als Bild der Intensität auf der Port-Ebene (der Ebene durch den Schwerpunkt des
        Lichts, senkrecht zu seiner mittleren Richtung) – Strahlen als kohärentes Strahlbild,
        Felder als Feldbild. Das Fenster umfasst das Licht mit einem Rand (Parameter
        \param{margin}).
\end{description}

Die Ergebnisse der Anzeigen stehen auch im Dock \gui{Results} (Ergebnissatz des Laufs), öffnen
aber nicht das allgemeine Image-Fenster oder Plot-Dock.

\subsection{Parameterstudien mit Meta-Blöcken}\label{sec:workflow-meta}

Meta-Blöcke wiederholen den Teil des Workflows vor einem Ausgang mit geänderten Parametern und
sammeln die Ergebnisse:

\begin{description}
  \item[\block{meta.sweep}] ein Parameter über eine Liste von Werten (Wellenlänge, Kippung,
        Defokus, Gangunterschied …),
  \item[\block{meta.replicates}] Wiederholungen mit Wiederholungswerten (z.\,B.\ Zufallszahlen),
        die gesammelten Bilder werden gemittelt,
  \item[\block{meta.optimise}] ein Parameter innerhalb von Grenzen, bis eine Messgröße des
        Ergebnisses ihr Maximum oder Minimum erreicht.
\end{description}

Den Parameter geben Sie als \code{Block.Parameter} oder \code{Block.Parameter[i]} an (z.\,B.\
\code{mirror.point[0]}), den gesammelten Ausgang als \code{Block.Anschluss} (z.\,B.\
\code{camera.image}). Stattdessen können Sie den Meta-Block auch verschalten: \code{drive} in
den \code{drive}-Eingang eines Spiegels, Abschnitts oder Strahlteilers (der Block verschiebt sich
längs seiner Flächennormalen um den Wert in mm – ein Piezo-Phasenschieber) oder in den
\code{pose}-Eingang eines Blocks (Komponente \param{component}), und den zu beobachtenden Ausgang in
\code{observe}. Die Messgrößen eines Bildes sind \code{peak} (Maximum), \code{centre} (Mittel der
$3 \times 3$ mittleren Pixel) und \code{power} (Summe). Verbinden Sie \code{results} mit einem
\block{display.plot}, um das Ergebnis zu sehen. Dank des Caches rechnet jede Wiederholung nur, was
hinter dem geänderten Block liegt.

% ---------------------------------------------------------------------------
\section{Gespeicherte Ergebnisse}\label{sec:workflow-gespeichert}

Beim Speichern schreibt die Anwendung das \emph{neueste Ergebnis jedes Anzeigefensters} –
Kamerabilder, Pupillenbilder, Kurven von Sweeps, Tabellen, Balken – mit in die
\datei{.rttf}-Datei (Abschnitt \code{results}, Anhang~\ref{kap:dateiformat}), jeweils mit dem
Fingerabdruck dessen, woraus es gerechnet wurde. Nach dem Öffnen

\begin{itemize}
  \item stehen diese Ergebnisse im Ergebnissatz \gui{Saved results} und in den Fenstern ihrer
        Anzeigeblöcke,
  \item gilt eine Anzeige, deren Fingerabdruck noch zum Workflow und seinen Einstellungen passt,
        als aktuell (grüner Haken), und \gui{Run workflow} rechnet sie nicht neu
        (\code{(unchanged)} in der Statuszeile),
  \item wird eine Anzeige neu gerechnet, sobald sich etwas geändert hat, wovon sie abhängt.
\end{itemize}

Die Zwischenergebnisse (Strahlbündel, Felder) werden \emph{nicht} gespeichert. Ein Lauf mit neuen
Ergebnissen markiert das Dokument als geändert.

% ---------------------------------------------------------------------------
\section{Eingebettete Dateien und Glaskataloge}\label{sec:workflow-eingebettet}\index{Glaskatalog}

Ein Dokument soll auf jedem Rechner rechnen – auch ohne die Dateien, aus denen es entstanden ist.
Deshalb speichern \gui{Check workflow} und \gui{Run workflow} Dateien, die Blöcke nennen,
\emph{in den Blöcken} und leeren die Pfade:

\begin{itemize}
  \item bei \block{fg.grating} die Jones-Tabelle (\param{table} $\to$ \param{table\_data}) und die
        FourierGrating-Konfiguration (\param{fouriergrating} $\to$ \param{config}),
  \item bei \block{fg.stack} die Konfiguration des Stapels.
\end{itemize}

Das Log nennt jede gespeicherte Datei. Kann eine genannte Datei nicht gelesen werden und hält der
Block keine Kopie, ist das eine Warnung (\code{graph.file\_unreadable}): Der Block rechnet erst,
wenn die Datei einmal hier lesbar war. STEP-Dateien von CAD-Schichten in einer
FourierGrating-Konfiguration bleiben außerhalb (\code{graph.file\_external}).

\subsection{Glaskataloge}

Raytatouille-Systeme brauchen für reale Gläser Glaskataloge (\datei{*.agf}). Der \emph{eine} Ort
dafür ist das Feld \gui{Glass catalogues} im Panel des Raytracers. Bei jedem Prüfen und Rechnen
übernimmt die Workflow-Maschine

\begin{itemize}
  \item die dort aufgeführten Dateien (eine Datei ohne Ordner wird neben der Systemdatei bzw.\
        im Ordner des Raytracer-Plugins gesucht) und
  \item alle \datei{*.agf} neben der Systemdatei und in deren Unterordner \datei{catalogs}
\end{itemize}

in die Quelle des Systems (\param{rtt\_catalogs}) und speichert ihren Inhalt im Dokument
(\param{rtt\_catalog\_data}). Danach rechnet die \datei{.rttf}-Datei auch ohne die
\datei{.agf}-Dateien – im Workflow wie im Raytracer-Panel, das bei fehlender Datei die
gespeicherten Kataloge verwendet. Ein Katalog, der sich nicht lesen lässt, hält den Lauf nicht
an: Log und Statuszeile nennen den Katalog, die fehlenden Gläser und die Abhilfe.

% ---------------------------------------------------------------------------
\section{Blöcke anderer Plugins}\label{sec:workflow-anbieter}

Der Workflow vermerkt für jeden Blocktyp, welches Plugin (Kennung und Version) ihn liefert.
Startet ein solches Plugin nicht oder ist es nicht installiert, bleibt der Block mit allen Daten im
Dokument; die Prüfung markiert ihn rot mit

\begin{lstlisting}
graph.provider_missing: block type 'fg.grating' comes from the plugin fg
(version 0.1.0), which is not available – the block keeps its data;
install or start the plugin
\end{lstlisting}

Installieren bzw.\ starten Sie das Plugin (\menu{Plugins,Reload plugins}) und prüfen Sie erneut.
Ein Plugin, das während einer Rechnung abstürzt, lässt nur seinen Block scheitern.

Ein Detektor, auf dessen Weg ein Block eines anderen Plugins liegt (z.\,B.\ eine Gitterprobe),
wird von der Workflow-Maschine selbst gerechnet, die beide Plugins erreicht.

\begin{hinweis}
Ein Workflow aus der Zeit, bevor die Blöcke ihre Daten selbst hielten, kann nicht gerechnet
werden; \gui{Run workflow} meldet \gui{This workflow holds no Raytatouille system data (made
before they were kept): run Workflow from system once in the Raytracer plugin, then Run workflow
again.} Laden Sie das System im Raytracer und wählen Sie einmal \gui{Workflow from system}.
\end{hinweis}

% ---------------------------------------------------------------------------
\section{Beispiel: ein neues System mit Kamera}\label{sec:workflow-beispiel}

Das folgende Beispiel baut einen kleinen Workflow von Grund auf. Es braucht den Raytracer und eine
eingerichtete Python-Umgebung.

\begin{schritte}
  \item Neues Dokument (\keys{Strg+N}). Im Plugin-Dock den \plugin{Raytracer} wählen und
        \gui{New Raytatouille system} klicken. Das Fenster \gui{Node workflow} öffnet sich mit
        einer Quelle, einer Blende, einem Detektor und dem Ast \code{main}.
  \item Doppelklick auf die Quelle: Wellenlängen, Feldpunkte und Apertur prüfen und ggf.\
        ändern.
  \item Rechtsklick auf die freie Fläche $\to$ \gui{Add block} $\to$ \code{domain} $\to$
        \block{rtt.part.lens}. Doppelklick auf die Linse und im Dialog Lage, Radien, Dicke und
        Glas eingeben.
  \item Die Verbindung zwischen Blende und Detektor löschen (Rechtsklick $\to$
        \gui{Delete connection}) und neu verbinden: Blende $\to$ \code{front} der Linse,
        \code{back} der Linse $\to$ Detektor.
  \item Den alten Ast \code{main} löschen und einen neuen aufnehmen
        (Abschnitt~\ref{sec:workflow-ast-aufnehmen}): Port der Quelle, Port, durch den das
        Licht die Blende verlässt, \code{back} der Linse, Detektor.
  \item \gui{Add block} $\to$ \code{analysis} $\to$ \block{display.image}. Den Ausgang
        \code{image} des Detektors mit dem Eingang \code{in} der Anzeige verbinden.
  \item Im Plugin \plugin{Workflow} \gui{Check workflow} klicken. Fehler sind rot markiert;
        das Log nennt sie.
  \item \gui{Run workflow} klicken. Danach Doppelklick auf die Anzeige: Das Fenster
        \gui{Display: image} zeigt das Bild.
  \item Einen Parameter der Linse ändern: Linse, Detektor und Anzeige werden rot. Ein neuer Lauf
        rechnet nur diese.
\end{schritte}

Parallel zeichnet der Raytracer das System im CAD-Modell und aktualisiert die Zeichnung nach
jeder Änderung.

% ---------------------------------------------------------------------------
\section{Meldungen und Fehler}\label{sec:workflow-fehler}

Die Prüfung meldet Probleme mit einem Code, einer Meldung und dem Ort (Block, Port, Verbindung
oder Ast). \emph{Fehler} (error) verhindern das Rechnen, \emph{Warnungen} (warning) nicht.

\begin{referenz}{@{}>{\raggedright\arraybackslash}p{0.4\linewidth}>{\raggedright\arraybackslash}X@{}}
\toprule
\textbf{Code} & \textbf{Bedeutung und Abhilfe} \\
\midrule
\endhead
\multicolumn{2}{@{}l}{\textit{Blöcke und Links}} \\
\code{graph.node\_type\_unknown} & Unbekannter Blocktyp. Plugin fehlt oder ist zu alt. \\
\code{graph.provider\_missing} & Der Typ stammt von einem Plugin, das nicht verfügbar ist
  (Abschnitt~\ref{sec:workflow-anbieter}). \\
\code{check.no\_evaluator} & Warnung: Für den Blocktyp gibt es in dieser Installation keinen
  Rechner (geplanter Typ, oder sein Plugin ist nicht geladen). Ein Lauf, der den Block braucht,
  bricht mit \code{exec.no\_evaluator} ab; Plugin laden oder Block entfernen. \\
\code{graph.param\_unknown} & Ein Parameter gehört nicht zum Blocktyp. \\
\code{graph.socket\_unknown} & Ein Link nennt einen Anschluss, den der Block nicht hat. \\
\code{graph.link\_direction} & Ein Link läuft nicht von einem Ausgang zu einem Eingang. \\
\code{graph.kind\_mismatch} & Licht-Anschluss mit Daten-Anschluss verbunden. \\
\code{graph.data\_type\_mismatch} & Die Datentypen der beiden Enden passen nicht. \\
\code{graph.no\_common\_representation} & Die beiden Enden haben keine gemeinsame
  Darstellung. \\
\code{graph.transport\_unsupported} & Der feste Transport wird nicht von beiden Enden
  angenommen. \\
\code{graph.input\_multiple} & Mehrere Links in einen Eingang; einen kombinierenden Block
  verwenden. \\
\code{graph.cycle} & Rückkopplungsschleife; dafür gibt es Meta-Blöcke. \\
\code{graph.detector\_unconnected} & Warnung: Detektor ohne Licht. \\
\multicolumn{2}{@{}l}{\textit{Ports, Verbindungen und Äste}} \\
\code{graph.port\_unknown} & Eine Verbindung nennt einen Port, den der Block nicht hat. \\
\code{graph.port\_links} & Mehrere Verbindungen an einem Port. \\
\code{graph.port\_facing} & Die verbundenen Ports zeigen nicht zueinander. \\
\code{graph.port\_behind} & Ein Port liegt hinter dem anderen. \\
\code{graph.port\_medium} & Außerhalb der beiden Ports liegen verschiedene Medien. \\
\code{graph.port\_aperture} & Warnung: Die Aperturen überlappen nicht. \\
\code{graph.port\_near\_field} & Warnung: Ports sehr nah beieinander; evaneszente Felder werden
  nicht übertragen. \\
\code{graph.port\_tilt} & Warnung: Felder zwischen gekippten Ebenen werden geometrisch
  umgetastet. \\
\code{graph.no\_branch} & Das Modell hat keinen Ast, oder an einem Detektor endet keiner. \\
\code{graph.branch\_duplicate} & Zwei Äste mit gleichem Namen oder gleichen Schritten. \\
\code{graph.branch\_impossible} & Ein Schritt ist so nicht möglich (z.\,B.\ ohne Verbindung,
  nicht von Quelle zu Detektor). \\
\code{graph.branch\_backend} & Kein Pfad des Systems passt zum Schritt durch einen
  Raytatouille-Block, oder ein Geisterbild passt nicht zu den Ports. \\
\code{graph.observer\_ambiguous} & Ein beobachteter Ausgang wird von mehreren Ästen durchlaufen;
  den Ast benennen. \\
\multicolumn{2}{@{}l}{\textit{Gruppen}} \\
\code{graph.group\_variant} & Die aktive Variante existiert nicht. \\
\code{graph.group\_ports}, \code{graph.group\_params} & Varianten exportieren
  unterschiedliche Ports bzw.\ Parameter. \\
\code{graph.group\_branch} & Ein Schritt nennt einen inneren Ast, den die Gruppe nicht hat. \\
\multicolumn{2}{@{}l}{\textit{Dateien}} \\
\code{graph.file\_unreadable} & Warnung: eine genannte Datei bzw.\ ein Glaskatalog ist nicht
  lesbar und nicht gespeichert. \\
\code{graph.file\_external} & Warnung: eine CAD-Schicht einer FourierGrating-Konfiguration
  bleibt eine externe Datei. \\
\multicolumn{2}{@{}l}{\textit{Planer}} \\
\code{plan.representation\_unavailable} & Der feste Transport kann noch nicht gerechnet werden
  (verfügbar z.\,B.\ nur R bzw.\ RH). \\
\code{plan.no\_representation} & Keine zulässige Darstellung für eine Verbindung. \\
\code{plan.no\_operator} & Keine Umwandlung zwischen den geplanten Darstellungen. \\
\code{plan.incoherent\_merge} & Warnung: Wege ohne Phase addieren sich als Intensitäten. \\
\code{plan.paths\_truncated} & Warnung: zu viele Wege; die Aufzählung wurde abgebrochen. \\
\code{plan.no\_output} & Warnung: kein Detektor und keine Auswertung, für die geplant werden
  kann. \\
\multicolumn{2}{@{}l}{\textit{Beim Rechnen}} \\
\code{exec.transport\_unsupported} & Der feste Transport kommt erst in einer späteren Stufe
  (G, T, P); meist fehlt \code{fieldtracing} für H bzw.\ S. \\
\code{exec.no\_evaluator} & Für den Blocktyp gibt es keinen Rechner (geplanter Typ, oder
  Plugin bzw.\ Paket fehlt); \gui{Check workflow} warnt davor mit \code{check.no\_evaluator}. \\
\code{exec.convert\_unsupported} & Ein \block{fusion.convert} soll in eine Darstellung umwandeln,
  die diese Version noch nicht rechnet (G, P), oder für R, H bzw.\ S fehlt \code{fieldtracing}. \\
\code{exec.spectrum\_cut} & Warnung: Ein Teil der Leistung fiel aus einem Spektrum (Anteil in der
  Meldung), meist evaneszente Wellen an einer Grenzfläche oder einer gekippten Ebene, oder
  Wellen, die einen Block nicht erreichen. Bei kleinen Anteilen harmlos; sonst die Apertur bzw.\
  die Ebene prüfen. Markiert die Verbindung bzw.\ den Block (Warnzeichen). \\
\code{exec.spectrum\_rays} & Warnung: Ein Spektrum wurde zu Strahlen, obwohl das Bündel nicht groß
  gegen die Wellenlänge ist oder sich Beugungsordnungen überlappen; die Strahlen verlieren die
  Interferenz. Abhilfe: die Verbindung bis zum Detektor auf \code{S} bzw.\ \code{H} lassen. \\
\code{exec.spectrum\_support} & Ein Spektrum ohne Abtastung (hinter einer verdichteten Gruppe)
  kann kein Feld und keine Strahlen werden. Die Meldung nennt den Block; die Verbindung in ihn
  auf \code{H} oder \code{R} stellen -- ein Detektor nimmt das Spektrum weiter als Summe. \\
\code{exec.param} & Ein Parameter ist ungültig, z.\,B.\ eine Datei nicht lesbar. \\
\code{exec.field} & Fehler bei einer Feldrechnung. \\
\code{exec.condense} & Die Gruppe lässt sich nicht verdichten. \\
\code{exec.cycle} & Zyklische Abhängigkeit, z.\,B.\ ein Meta-Block sammelt in sich selbst. \\
\bottomrule
\end{referenz}

Weitere typische Meldungen von \gui{Run workflow}:

\begin{referenz}{@{}>{\raggedright\arraybackslash}p{0.42\linewidth}>{\raggedright\arraybackslash}X@{}}
\toprule
\textbf{Meldung} & \textbf{Abhilfe} \\
\midrule
\endhead
\code{No node workflow yet – …} & Erst einen Workflow anlegen
  (Abschnitt~\ref{sec:workflow-entstehen}). \\
\code{Nothing to show: attach a display block …} & Einen Anzeigeblock an eine Kamera, einen
  Phasenschieber oder einen Lichtausgang hängen. \\
\code{The workflow has errors – Check workflow marks them: …} & \gui{Check workflow} ausführen
  und die rot markierten Stellen beheben. \\
\code{This workflow holds no Raytatouille system data …} & Einmal \gui{Workflow from system}
  im Raytracer. \\
\code{The node workflow needs the Python package rtt\_fusion …} & Paket \code{rtt-fusion} in
  die Plugin-Umgebung installieren (Abschnitt~\ref{sec:umgebung-pruefen}). \\
\code{workflow: nothing computed – …} & Den genannten Grund beheben; weitere stehen im Log. \\
\textit{Anzeige}\code{: dark – the light focuses into less than a pixel; …} & Wellenoptik
  einschalten oder Verbindung auf H und kleines Bildfenster. \\
\bottomrule
\end{referenz}
